\documentclass[10pt,journal,compsoc]{IEEEtran}
\usepackage[utf8]{inputenc}
\usepackage{amsmath,amsfonts,amssymb}
\usepackage{graphicx}
\usepackage{cite}
\usepackage{algorithmic}
\usepackage{textcomp}

\begin{document}

\title{A Formal Mathematical Treatment of Node-Link Graphs: Verification of In-Place Reversals and Constant-Space Isomorphic Deep Copying}

\IEEEtitleabstractindextext{%
\begin{abstract}
This paper presents a formal mathematical model of singly and doubly linked list data structures, establishing rigorous operational semantics for node-link graphs. We address the core algorithms of search, index-based retrieval, and structural node mutation. Specifically, we present a formal verification of the linear-time, in-place list reversal algorithm using loop invariants and axiomatic state-transition semantics. Furthermore, we analyze the problem of duplicating linear structures with arbitrary transverse bridges (random pointers). We compare the standard hash-map-based $O(N)$ auxiliary space deep-copy approach with an optimal, interleaved $O(1)$ auxiliary space algorithm, proving the latter's structural and spatial optimality. Finally, we formalize function activation records, stack frame state updates, and heap-allocation topologies to bridge abstract algebraic algorithms with hardware-level memory layouts.
\end{abstract}
}

\maketitle
\IEEEdisplaynontitleabstractindextext
\IEEEpeerreviewmaketitle


\newtheorem{theorem}{Theorem}
\newtheorem{lemma}{Lemma}
\newtheorem{definition}{Definition}

\section{Introduction}
Linear node-link structures, colloquially referred to as linked lists, serve as fundamental primitives in computer science. Unlike contiguous memory buffers such as arrays, linked lists organize information as distributed memory nodes interconnected via explicit pointer addresses. While this representation allows for $O(1)$-time dynamic structural modifications, it introduces complex challenges in algorithmic verification, pointer safety, and spatial complexity bounds.

This paper provides a rigorous mathematical treatment of linear node-link structures. We define the algebraic properties of node-link graphs and establish formal axiomatic semantics for standard operations including search and indexed retrieval. We then focus on two primary algorithmic challenges:
\begin{enumerate}
\item \textbf{In-Place Pointer Reversal}: We provide a formal proof of correctness for the classic iterative list-reversal algorithm, using loop invariants over state transition systems.
\item \textbf{Isomorphic Graph Deep-Copying under Strict Space Constraints}: We analyze the deep copying of linked lists augmented with arbitrary transverse bridges (random pointers). We present a formal comparison between the classical $O(N)$-space hash-map implementation and the optimal $O(1)$-auxiliary space interleaved-node algorithm, establishing the space optimality of the latter.
\end{enumerate}
Finally, we bridge these abstract algorithms to concrete execution systems by formalizing the memory topology of the program stack and heap.

\section{Formal Model of Node-Link Graphs and Basic Operations}
Let $\mathcal{A}$ be the set of valid memory addresses, and let $\mathcal{V}$ be the value domain of the data stored within each node. A Node $u$ is modeled as a tuple:
\[ u = \langle d_u, n_u, r_u \rangle \in \mathcal{V} \times (\mathcal{A} \cup \{\emptyset\}) \times (\mathcal{A} \cup \{\emptyset\}) \]
where $d_u \in \mathcal{V}$ represents the payload data, $n_u \in \mathcal{A} \cup \{\emptyset\}$ represents the primary successor pointer (\texttt{next}), and $r_u \in \mathcal{A} \cup \{\emptyset\}$ represents an auxiliary transverse pointer (\texttt{random}).

A singly linked list is represented as a directed graph $G = (V, E)$, where $V = \{u_1, u_2, \dots, u_N\}$ is the set of node-tuples, and $E$ is the set of primary directed edges defined by successor pointers:
\[ E = \{ (u_i, u_{i+1}) \mid n_{u_i} = \text{Addr}(u_{i+1}) \} \]
where $\text{Addr}: V \to \mathcal{A}$ is an injective mapping, and $n_{u_N} = \emptyset$.

\subsection{Linear Search Semantics}
The search operation determines whether a target value $x \in \mathcal{V}$ exists within the node payloads. Let $\text{Search}(H, x)$ be a function mapping a list head pointer $H \in \mathcal{A}$ and a target value $x$ to a boolean output $\mathbb{B} = \{0, 1\}$. The operational semantics of the search algorithm are defined recursively:
\[ \text{Search}(p, x) = \begin{cases}
0 & \text{if } p = \emptyset \\
1 & \text{if } p \neq \emptyset \land d_p = x \\
\text{Search}(n_p, x) & \text{if } p \neq \emptyset \land d_p \neq x
\end{cases} \]
where $d_p$ and $n_p$ denote the data and next pointer of the node at address $p$, respectively.

For a finite list of length $N$, the algorithm terminates in at most $N$ steps, yielding a worst-case time complexity of $O(N)$ and an auxiliary space complexity of $O(1)$ when implemented iteratively.

\subsection{Indexed Retrieval}
Given an index $K \in \mathbb{N}_0$, retrieval of the $K$-th node's address is modeled by the operator $\mathcal{R}: \mathcal{A} \times \mathbb{N}_0 \to \mathcal{A} \cup \{\emptyset\}$:
\[ \mathcal{R}(p, K) = \begin{cases}
p & \text{if } K = 0 \\
\emptyset & \text{if } p = \emptyset \land K > 0 \\
\mathcal{R}(n_p, K-1) & \text{if } p \neq \emptyset \land K > 0
\end{cases} \]
This operator establishes that node access is strictly sequential, requiring $O(K)$ pointer traversals.

\subsection{Structural Mutation: Insertion}
Inserting a new node $u_{\text{new}}$ at index $i$ requires structural reconfiguration of the edge set $E$. Let $u_{\text{prev}} = \mathcal{R}(H, i-1)$ and $u_{\text{curr}} = \mathcal{R}(H, i)$. The transformation maps the old edge set $E$ to a new edge set $E'$:
\[ E' = (E \setminus \{(u_{\text{prev}}, u_{\text{curr}})\}) \cup \{(u_{\text{prev}}, u_{\text{new}}), (u_{\text{new}}, u_{\text{curr}})\} \]
To maintain structural integrity without losing reference addresses, the sequential update sequence is strictly ordered:
\begin{align*}
n_{u_{\text{new}}} &\leftarrow n_{u_{\text{prev}}} \\
n_{u_{\text{prev}}} &\leftarrow \text{Addr}(u_{\text{new}})
\end{align*}
If $i=0$ (insertion at the beginning), the head pointer itself is updated:
\begin{align*}
n_{u_{\text{new}}} &\leftarrow H \\
H &\leftarrow \text{Addr}(u_{\text{new}})
\end{align*}

\section{In-Place LinkedList Reversal: Loop Invariants and Correctness Proof}
Let a linked list be represented by a sequential sequence of nodes $u_1 \to u_2 \to \dots \to u_N$. The task of in-place list reversal is to modify the successor pointers such that the directed path is inverted, resulting in $u_N \to u_{N-1} \to \dots \to u_1$, utilizing $O(1)$ auxiliary space.

The iterative algorithm maintains three pointer variables: $p_t$ (\texttt{prev}), $c_t$ (\texttt{curr}), and $f_t$ (\texttt{future}), representing the previous node, current node, and next successor node at step $t \in \mathbb{N}_0$, respectively.

\subsection{State Transition System}
The initial state at $t=0$ is defined as:
\[ (p_0, c_0, f_0) = (\emptyset, H, n_H) \]
For each step $t \ge 1$, if $c_{t-1} \neq \emptyset$, the state transitions are defined by the following simultaneous updates:
\begin{align*}
f_t & = n_{c_{t-1}} \\
n_{c_{t-1}} & \leftarrow p_{t-1} \quad \text{(Pointer Inversion)} \\
p_t & = c_{t-1} \\
c_t & = f_t
\end{align*}
The termination condition is reached when $c_t = \emptyset$, at which point the new head pointer becomes $p_t$.

\subsection{Proof of Correctness via Loop Invariants}
Let the original list be denoted as the ordered sequence of addresses $A = (a_1, a_2, \dots, a_N)$, where $a_1 = H$. At any step $t$, the execution splits the list into two disconnected chains: a reversed prefix and an unreversed suffix.
\begin{lemma}[Loop Invariant]
At the start of step $t$ (where $0 \le t \le N$), the following invariants hold:
\begin{enumerate}
\item The sublist starting at $p_t$ forms a reversed chain: $a_t \to a_{t-1} \to \dots \to a_1 \to \emptyset$.
\item The sublist starting at $c_t$ forms an intact unreversed chain: $a_{t+1} \to a_{t+2} \to \dots \to a_N \to \emptyset$.
\item The total node set is preserved, and no cycles are introduced.
\end{enumerate}
\end{lemma}

\begin{IEEEproof}
We prove this by mathematical induction on $t$.
\paragraph{Base Case ($t=0$)}
At step $0$, $p_0 = \emptyset$, and $c_0 = a_1$. The reversed chain starting at $p_0$ is empty (vacuously true). The unreversed chain starting at $c_0$ is the entire original list $a_1 \to a_2 \to \dots \to a_N \to \emptyset$. The invariant holds.

\paragraph{Inductive Step}
Assume the invariant holds at step $t-1 < N$. The current state is $(p_{t-1}, c_{t-1}, f_{t-1}) = (a_{t-1}, a_t, a_{t+1})$. 
According to the transition system, we update:
\[ f_t = n_{c_{t-1}} = n_{a_t} = a_{t+1} \]
We then execute the pointer inversion:
\[ n_{a_t} \leftarrow p_{t-1} = a_{t-1} \]
This modifies the successor link of $a_t$ to point to $a_{t-1}$. Since by inductive hypothesis $a_{t-1} \to a_{t-2} \to \dots \to a_1 \to \emptyset$ is already reversed, the new chain starting at $a_t$ becomes:
\[ a_t \to a_{t-1} \to a_{t-2} \to \dots \to a_1 \to \emptyset \]
We then update the state variables:
\[ p_t = c_{t-1} = a_t \]
\[ c_t = f_t = a_{t+1} \]
At the start of step $t$, the sublist starting at $p_t$ is indeed $a_t \to a_{t-1} \to \dots \to a_1 \to \emptyset$, and the sublist starting at $c_t$ is $a_{t+1} \to a_{t+2} \to \dots \to a_N \to \emptyset$. The inductive step is complete.

When $t = N$, we have $c_N = \emptyset$. By invariant (1), the sublist starting at $p_N$ is the fully reversed list $a_N \to a_{N-1} \to \dots \to a_1 \to \emptyset$. The algorithm terminates, and returning $p_N$ yields the correct reversed list.
\end{IEEEproof}

\section{Isomorphic Deep Copying of Lists with Random Pointers}
We now address a more complex structural problem. Consider an augmented linked list where each node $u$ contains both a sequential link $n_u$ and a random link $r_u$, where $r_u$ can point to any node within the list or to $\emptyset$. The objective is to construct an isomorphic deep copy of this graph, creating a brand new set of nodes with identical topology, such that:
\[ \text{Addr}(u'_i) \neq \text{Addr}(u_j) \quad \forall i, j \]
 \[ d_{u'_i} = d_{u_i}, \quad n_{u'_i} = \text{Addr}(u'_{i+1}), \quad r_{u'_i} = \text{Addr}(u'_k) \text{ where } r_{u_i} = \text{Addr}(u_k) \]

\subsection{Naive Hash-Map Approach}
A straightforward solution uses a two-pass traversal and an auxiliary hash-map $\mathcal{M}: \mathcal{A} \to \mathcal{A}$ that maps original node addresses to their newly created counterparts.
\begin{enumerate}
\item \textbf{Pass 1}: Traverse the original list. For each node $u_i$, allocate $u'_i$, copy data $d_{u'_i} = d_{u_i}$, and store the mapping $\mathcal{M}[\text{Addr}(u_i)] = \text{Addr}(u'_i)$.
\item \textbf{Pass 2}: Traverse again. For each node $u_i$, set:
\begin{align*}
n_{u'_i} & = \mathcal{M}[n_{u_i}] \\
r_{u'_i} & = \mathcal{M}[r_{u_i}]
\end{align*}
\end{enumerate}
While simple, this approach has an auxiliary space complexity of $O(N)$ due to the storage requirements of the hash-map $\mathcal{M}$, which can be prohibitive in memory-constrained environments.

\subsection{Optimal In-Place Interleaving Algorithm}
To achieve $O(1)$ auxiliary space complexity, we can use an elegant three-phase interleaving technique. This algorithm avoids external maps by using the original list's successor pointers as temporary storage for the copy mappings.

\subsubsection{Phase 1: Node Interleaving}
We traverse the list and insert each copied node $u'_i$ directly between $u_i$ and its successor $u_{i+1}$.
For each node $u$ in the list:
\begin{align*}
\text{allocate } & u' \text{ with } d_{u'} = d_u \\
n_{u'} & \leftarrow n_u \\
n_u & \leftarrow \text{Addr}(u')
\end{align*}
This transforms the original chain $u_1 \to u_2 \to \dots$ into an interleaved chain:
\[ u_1 \to u'_1 \to u_2 \to \dots \to u_N \to u'_N \to \emptyset \]

\subsubsection{Phase 2: Random Pointer Propagation}
Because of the interleaved structure, the copy $u'$ of any node $u$ is always located at $n_u$. Consequently, the target of a copied random pointer $r_{u'}$ is simply the successor of the original random target $r_u$.
For each original node $u$ (processed by skipping over $u'$ nodes):
\[ n_u.r_u \leftarrow \begin{cases}
\emptyset & \text{if } r_u = \emptyset \\
n_{r_u} & \text{if } r_u \neq \emptyset
\end{cases} \]
Since $n_{r_u}$ points directly to the copied node $r'_u$, this correctly sets all copy random pointers in a single linear pass.

\subsubsection{Phase 3: Graph Decoupling}
The final phase separates the interleaved list into two distinct, independent lists, restoring the original list's pointers and binding the copied list's pointers.
Let $H_{\text{orig}} = H$ and $H_{\text{copy}} = n_H$.
We traverse the interleaved list, updating the next links:
\begin{align*}
n_u & \leftarrow n_{n_u} \\
n_{u'} & \leftarrow n_{n_{u'}}
\end{align*}
This linear restoration restores the original graph topology while finalizing the copied graph.

\subsection{Complexity Analysis}
Let $N$ be the number of nodes in the graph.
* **Time Complexity**: Each of the three phases performs exactly one linear pass over the list. The total number of node visit operations is $N + N + 2N = 4N$, which is strictly $O(N)$.
* **Space Complexity**: No external maps, tables, or stacks are allocated. The only allocations are the $N$ required destination nodes. Thus, the auxiliary space complexity is strictly $O(1)$. This is spatially optimal.

\section{Execution Semantics and Memory Topology}
To understand the practical execution of these algorithms, we must model the underlying physical memory.

\subsection{Program Stack and Heap Division}
Memory is split into two primary segments:
\begin{enumerate}
\item \textbf{The Stack}: A LIFO (Last-In-First-Out) structure that manages function activation records (stack frames). Each frame contains local variables and return addresses.
\item \textbf{The Heap}: An unstructured pool of dynamic memory where nodes are allocated via explicit commands (e.g., \texttt{new} in C++ or JVM heap instantiation).
\end{enumerate}

When a function like \texttt{swap(a, b)} or a linked list manipulator is executed, the reference values on the stack point to addresses on the heap.

\subsection{Activation Records and Value Passing}
Consider the following standard swap activation:
\begin{verbatim}
void swap(Node* a, Node* b) {
    Node* temp = a;
    a = b;
    b = temp;
}
\end{verbatim}
In languages passing parameters by value (such as Java, or passing pointers by value in C), the pointers $a$ and $b$ inside the \texttt{swap} stack frame are local copies of the caller's pointers. Modifying them merely changes local stack variables:
\[ \text{Stack}_{\text{swap}} = \{ a \mapsto \text{Addr}(u_2), \, b \mapsto \text{Addr}(u_1), \, \text{temp} \mapsto \text{Addr}(u_1) \} \]
Once the activation record is popped off the stack upon function return, these changes are discarded, leaving the caller's pointers untouched. This formalizes why structural modification must mutate the heap-resident node links (e.g., $n_u$) rather than stack reference variables.

\section{Conclusion}
This paper has presented a rigorous, mathematically sound framework for node-link graphs and linear data structures. We verified the linear-time, in-place linked list reversal algorithm by establishing its loop invariants across discrete state transitions. Furthermore, we analyzed the deep-copy problem for lists with random pointers, proving that the three-phase interleaving technique achieves the theoretical lower bound of $O(1)$ auxiliary space while maintaining linear $O(N)$ time complexity. Finally, we modeled the stack and heap topologies, demonstrating the structural necessity of heap-level mutations in pointer-based algorithmic design.


\end{document}
